%------------------------------------------------------------------------------%
\begin{question}
  Dada a equação da hipérbole
  \[
    \frac{(x-1)^2}{16} - \frac{(y+2)^2}{9} = 1
  \]
  \vspace{-1.5\baselineskip}
  \begin{enumerate}
    \item encontre o centro e os vértices
    \item escreva a equação matricial
    \item esboce a curva
  \end{enumerate}
\end{question}

%------------------------------------------------------------------------------%
\begin{resolution}{Centro}
  Comparando
  \[
    \frac{(x-1)^2}{16} - \frac{(y+2)^2}{9} = 1
  \]
  \pause
  com a equação transladada
  \[
    \frac{(x-h)^2}{a^2} - \frac{(y-k)^2}{b^2} = 1
  \]
  \pause
  identificamos
  \[
    \pause
    a^2 = 16
    \qquad
    \pause
    a = 4
    \qquad\qquad\qquad
    \pause
    b^2 = 9
    \qquad
    \pause
    b = 3
  \]
  \pause
  Centro: \quad
  \(
    \pause
    C
    = (h, k)
    \pause
    = (1, -2)
  \)
\end{resolution}

%------------------------------------------------------------------------------%
\begin{resolution}{Vértices}
\begin{columns}[t]
\column{0.6\textwidth}
  \onslide<+->{Novas coordenadas}
  \[
    \onslide<+->{V
    = \begin{bmatrix}
      x - h  \\
      y - k
    \end{bmatrix}}
    \onslide<+->{= \begin{bmatrix}[c]
      x - 1 \\
      y - (-2)
    \end{bmatrix}}
    \onslide<+->{= \begin{bmatrix}
      x - 1  \\
      y + 2
    \end{bmatrix}}
  \]
\column{0.4\textwidth}
  \onslide<+->{Condição nos vértices}
  \[
    \onslide<+->{V
    = \begin{bmatrix}
      x - 1  \\
      y + 2
    \end{bmatrix}
    = \begin{bmatrix}
      \pm 4  \\
      0
    \end{bmatrix}}
  \]
\end{columns}

\medskip
\onslide<+->{Resolvendo a condição nos vértices}
\vspace{-1.5\baselineskip}
\begin{columns}[t]
\column{0.2\textwidth}
  \begin{gather*}
    \onslide<+->{y + 2 = 0 \\}
    \onslide<+->{y     = - 2 }
  \end{gather*}
\column{0.2\textwidth}
  \begin{gather*}
    \onslide<+->{x - 1 = 4     \\}
    \onslide<+->{x     = 4 + 1 \\}
    \onslide<+->{x     = 5}
  \end{gather*}
\column{0.2\textwidth}
  \begin{gather*}
    \onslide<+->{x - 1 = -4 \\}
    \onslide<+->{x     = -4 + 1 \\}
    \onslide<+->{x     = -3}
  \end{gather*}
\end{columns}

  \medskip
  \onslide<+->{Vértices:}
  \(
    \quad
    \onslide<+->{\mathcal{V}_1 = (-3, -2)}
    \quad
    \onslide<+->{\mathcal{V}_1 = (5, -2)}
  \)
\end{resolution}

%------------------------------------------------------------------------------%
\begin{resolution}{Equação Matricial}
  Equação Matricial
  \[
    \pause
    V^t A V = 1
  \]
  \pause
  com
  \[
    \pause
    A
    = \begin{bmatrix}
      \dfrac{1}{16} & 0            \\[2mm]
      0             & -\dfrac{1}{9}
    \end{bmatrix}
    \qquad\qquad
    \pause
    V = \begin{bmatrix}
      x - 2  \\
      y + 1
    \end{bmatrix}
  \]
\end{resolution}

%------------------------------------------------------------------------------%
\begin{resolution}{Esboço}
  \centering
  \includegraphics{I-06-Hiperbole-ex-3}
\end{resolution}


%------------------------------------------------------------------------------%
\endinput

